\section{Технический проект}

\subsection{Архитектура программной системы}

Разрабатываемая информационно-аналитическая система построена по архитектуре «клиент-сервер» с использованием REST-подхода.

Серверная часть реализована на языке Python с использованием фреймворка FastAPI и обеспечивает обработку запросов, выполнение бизнес-логики, а также взаимодействие с базой данных PostgreSQL и файловым хранилищем документов.

Клиентская часть представляет собой веб-приложение, доступ к которому осуществляется через браузер. Взаимодействие между клиентом и сервером выполняется по протоколу HTTP с использованием формата JSON для обмена данными.

Архитектура системы включает следующие основные компоненты:

\begin{itemize}
	\item клиентское веб-приложение (интерфейс пользователя);
	\item серверное приложение (FastAPI);
	\item система управления базами данных PostgreSQL;
	\item файловое хранилище для загружаемых и сгенерированных документов;
	\item модуль анализа и проверки документов формата DOCX.
\end{itemize}

Логика работы системы организована следующим образом. Пользователь через веб-интерфейс выполняет действия (авторизация, загрузка документа, запуск проверки и др.), которые инициируют HTTP-запросы к серверу. Сервер обрабатывает запрос, при необходимости обращается к базе данных и файловой системе, выполняет проверку документа и возвращает результат клиенту.

Модуль проверки документов осуществляет анализ структуры и оформления файлов формата DOCX, выявляет несоответствия установленным требованиям и формирует отчёт об ошибках. При необходимости система генерирует исправленный файл с визуальным выделением нарушений.

Все загружаемые документы сохраняются в файловой системе, а информация о них (путь к файлу, результаты проверки, пользователь и др.) — в базе данных PostgreSQL. Это обеспечивает разделение данных и повышает масштабируемость системы.

Для обеспечения безопасности и управления доступом используется механизм аутентификации пользователей. После успешного входа пользователь получает токен доступа, который используется при последующих запросах к API.

Общая схема взаимодействия компонентов системы представлена на рисунке~\ref{fig:architecture}.

\begin{figure}[H]
	\centering
	\includegraphics[width=1\textwidth]{diagramComponent.png}
	\caption{Общая схема архитектуры системы}
	\label{fig:architecture}
\end{figure}

\subsection{Модуль обработки учебно-методических документов}

Модуль обработки учебно-методических документов является центральным компонентом системы, отвечающим за приём, анализ и проверку файлов формата DOCX. Он одинаково работает для методических указаний, учебных пособий и монографий, поскольку алгоритм проверки оформления и механизмы работы с документами универсальны. Различия между типами документов проявляются только на этапе генерации титульных листов, где используются разные наборы полей и шаблоны оформления.

\subsubsection{Общая архитектура модуля}

Модуль построен по многоуровневой архитектуре и включает следующие компоненты:

\begin{itemize}
	\item \textbf{слой маршрутизации} — приём HTTP-запросов, валидация входных данных, маршрутизация к соответствующим обработчикам;
	\item \textbf{сервисный слой} — реализация бизнес-логики: проверка документов, генерация титульных листов, расчёт статистики;
	\item \textbf{слой доступа к данным} — объектно-реляционное отображение для взаимодействия с PostgreSQL;
	\item \textbf{слой утилит} — вспомогательные функции для работы с файлами, хеширования, логирования.
\end{itemize}

Такое разделение обеспечивает слабую связанность компонентов и упрощает тестирование и дальнейшее расширение.

\subsubsection{Интеграция с Microsoft Word через COM-интерфейс}

Ключевой особенностью модуля является использование COM-интерфейса Microsoft Word через библиотеку \texttt{pywin32}. При запуске проверки создаётся экземпляр приложения Word в фоновом режиме с отключёнными всплывающими окнами и диалогами. Это позволяет выполнять анализ и модификацию документов без видимого запуска графического интерфейса Word.

Основные операции, выполняемые через COM-интерфейс:

\begin{itemize}
	\item открытие документа (\texttt{Documents.Open}) с возможностью только чтения или с правом записи;
	\item получение общей информации о документе: количество страниц, разделов, абзацев;
	\item обход всех абзацев документа и извлечение параметров форматирования;
	\item доступ к структуре страниц (поля, колонтитулы, нумерация);
	\item модификация документа: изменение цвета шрифта для выделения ошибок, добавление комментариев;
	\item сохранение изменённой копии и закрытие документа.
\end{itemize}

\subsubsection{Алгоритм автоматической проверки оформления}

Процесс проверки запускается после получения файла и его временного сохранения. Алгоритм реализован в классе \texttt{WordMethodicalChecker} и включает следующие этапы.

\textbf{Этап 1. Предварительная обработка.}

Файл сохраняется во временную директорию, создаётся уникальный идентификатор задачи для отслеживания прогресса. Если требуется сохранение исправленной версии (параметр \texttt{mark\_document = True}), создаётся копия документа с суффиксом \texttt{\_checked}, с которой будет работать модуль. Исходный файл остаётся неизменным.

\textbf{Этап 2. Определение объёма работы.}

Через COM-интерфейс открывается документ, вычисляется общее количество страниц с помощью метода \texttt{ComputeStatistics(WD\_STATISTIC\_PAGES)}. Эта информация используется для последующей группировки ошибок по страницам.

\textbf{Этап 3. Фильтрация служебных элементов.}

При обходе абзацев пропускаются:

\begin{itemize}
	\item пустые абзацы — определяются по длине текста (менее 3 символов);
	\item абзацы, находящиеся внутри таблиц — метод \texttt{p.Range.Information(WD\_WITHIN\_TABLE)} возвращает \texttt{True} для таких элементов;
	\item заголовки — определяются по имени стиля: если в названии стиля присутствует подстрока \texttt{"heading"} или \texttt{"заголов"}, абзац считается заголовком и пропускается.
\end{itemize}

Отфильтрованные элементы не участвуют в сравнении параметров оформления, так как к ним предъявляются иные требования или они не являются частью основного текста.

\textbf{Этап 4. Извлечение параметров форматирования.}

Для каждого значимого абзаца извлекаются следующие характеристики:

\begin{itemize}
	\item \textbf{шрифт} — \texttt{p.Range.Font.Name}, эталонное значение "Times New Roman";
	\item \textbf{размер шрифта} — \texttt{p.Range.Font.Size}, эталонное значение 16 pt;
	\item \textbf{отступ первой строки} — \texttt{p.Range.ParagraphFormat.FirstLineIndent}, сначала проверяется формат самого абзаца, если значение близко к нулю, выполняется попытка получить отступ из стиля абзаца;
	\item \textbf{межстрочный интервал} — \texttt{p.Range.ParagraphFormat.LineSpacingRule}, эталонное значение \texttt{WD\_LINE\_SPACE\_SINGLE} (одинарный интервал);
	\item \textbf{выравнивание текста} — анализируется при проверке номеров страниц.
\end{itemize}

Особое внимание уделяется случаю, когда в одном абзаце встречается неоднородное форматирование (например, разные размеры шрифта в одной строке). В таких ситуациях выполняется дополнительный обход всех слов внутри абзаца через коллекцию \texttt{p.Range.Words}. Для каждого слова извлекается его размер шрифта \texttt{w.Font.Size}. Если размер отличается от эталонного значения 16 pt более чем на 0,2 pt, и при этом слово не является знаком препинания (точка, запятая, скобка и т.п.), фиксируется ошибка «неверный размер шрифта» с указанием конкретного слова и номера страницы. Значения размера, равные 9999999,0, интерпретируются как ошибочные (такое значение возникает при попытке прочитать размер шрифта у объекта, не имеющего единого форматирования) и также приводят к фиксации нарушения. Такой подход позволяет выявить ситуации, когда большая часть абзаца отформатирована правильно, но отдельные слова имеют неверный кегль, что часто встречается при копировании текста из внешних источников.

\textbf{Этап 5. Сравнение с эталонами и фиксация ошибок.}

Полученные параметры преобразуются в сантиметры (для отступов) и сравниваются с эталонными значениями:

\begin{itemize}
	\item для отступа используется допуск 0,05 см: \texttt{abs(indent\_cm - 1.25) > 0.05};
	\item для размера шрифта используется допуск 0,2 pt: \texttt{abs(size - 16) > 0.2};
	\item для шрифта и интервала выполняется точное сравнение (без допуска).
\end{itemize}

При обнаружении отклонения фиксируется номер страницы, тип нарушения и поясняющее сообщение. Номер страницы получается методом \texttt{p.Range.Information(WD\_ACTIVE\_END\_PAGE\_NUMBER)}. Каждая ошибка сохраняется в виде объекта \texttt{CheckIssue}, содержащего:

\begin{itemize}
	\item \texttt{rule} — раздел проверки (например, "3" — для параметров текста, "4" — для нумерации);
	\item \texttt{severity} — степень важности ("ERROR" или "WARNING");
	\item \texttt{location} — местоположение в документе (страница или диапазон страниц);
	\item \texttt{message} — текст сообщения с описанием нарушения и рекомендацией по исправлению;
	\item \texttt{priority} — приоритет (используется для сортировки).
\end{itemize}

\textbf{Этап 6. Проверка полей страницы.}

После обработки всех абзацев выполняется проверка параметров страниц. Для каждого раздела документа (коллекция \texttt{doc.Sections}) извлекаются и проверяются настройки \texttt{PageSetup}:

\begin{itemize}
	\item верхнее поле (\texttt{TopMargin}) — эталон 3,0 см;
	\item нижнее поле (\texttt{BottomMargin}) — эталон 2,25 см;
	\item левое поле (\texttt{LeftMargin}) — эталон 2,2 см;
	\item правое поле (\texttt{RightMargin}) — эталон 2,3 см.
\end{itemize}

Все поля проверяются в сантиметрах с допуском 0,05 см. При нарушении фиксируется номер раздела, и ошибка добавляется в отчёт.

\textbf{Этап 7. Проверка нумерации страниц.}

Проверка выполняется по всем разделам документа. Для каждого раздела анализируются верхний (\texttt{Headers(WD\_HEADER\_FOOTER\_PRIMARY)}) и нижний (\texttt{Footers(WD\_HEADER\_FOOTER\_PRIMARY)}) колонтитулы:

\begin{itemize}
	\item Если в верхнем колонтитуле обнаружены поля с типом \texttt{WD\_FIELD\_PAGE} — нумерация присутствует в верхней части страницы, дополнительно проверяется выравнивание номера (должно быть по центру, \texttt{WD\_ALIGN\_CENTER}).
	\item Если номер найден только в нижнем колонтитуле — фиксируется ошибка о неверном расположении номера страницы.
	\item Если номер не найден ни в верхнем, ни в нижнем колонтитуле — фиксируется ошибка об отсутствии нумерации страниц.
\end{itemize}

\textbf{Этап 8. Группировка ошибок.}

Для повышения удобочитаемости отчёта последовательные страницы с одинаковыми типами ошибок объединяются в диапазоны. Например, если ошибки шрифта обнаружены на страницах 5, 6, 7, 8, отчёт покажет «Стр. 5–8» вместо четырёх отдельных записей. Группировка выполняется методом \texttt{\_group\_pages}, который сортирует номера страниц, находит в них непрерывные последовательности и формирует соответствующие записи. Для каждого диапазона создаётся одна запись в отчёте, а для всех страниц диапазона выполняется визуальное выделение ошибок.

\textbf{Этап 9. Завершение проверки и освобождение ресурсов.}

По окончании обхода всех абзацев и проверки полей и нумерации страниц выполняется сохранение изменений в документе (если \texttt{mark\_document = True}) и закрытие документа через \texttt{doc.Close(False)}. После этого экземпляр Word завершается вызовом \texttt{self.word.Quit(False)}. Все временные файлы, созданные в процессе проверки (временная копия документа, файлы для отчёта), удаляются. Если на любом из этапов произошла ошибка, Word принудительно завершается, а временные файлы удаляются в блоке \texttt{finally}, что гарантирует отсутствие «висящих» процессов и освобождение дискового пространства.

\subsubsection{Формирование исправленной копии документа}

Если в документе обнаружены ошибки и установлен флаг \texttt{mark\_document} (активируется по умолчанию), модуль создаёт копию, в которой проблемные фрагменты визуально выделены. Алгоритм выделения работает следующим образом:

\begin{enumerate}
	\item Для каждой группы ошибок определяется диапазон страниц, на которых они обнаружены.
	\item Для первой страницы диапазона вычисляется начальная позиция через \texttt{doc.GoTo(1, 1, page)}. Константа 1 обозначает переход к странице, второй параметр — тип объекта (страница).
	\item Для последней страницы диапазона вычисляется конечная позиция через \texttt{doc.GoTo(1, 1, page + 1)}. Если страница последняя в документе, конечная позиция определяется как \texttt{doc.Content.End}.
	\item Область выделения задаётся через \texttt{doc.Range(start, end)}.
	\item Цвет шрифта в выделенной области устанавливается в красный (\texttt{WD\_COLOR\_RED}) с помощью свойства \texttt{rng.Font.Color}.
	\item К выделенной области добавляется комментарий с поясняющим сообщением через \texttt{doc.Comments.Add(rng, message)}.
\end{enumerate}

Такой подход позволяет сохранить исходную структуру документа и одновременно наглядно показать автору все места, требующие исправления. Комментарии Word видны при наведении курсора и не нарушают вёрстку документа.

\subsubsection{Генерация PDF-отчёта}

Параллельно с модификацией документа формируется PDF-отчёт с использованием библиотеки \texttt{reportlab}. Процесс создания отчёта включает следующие шаги:

\begin{enumerate}
	\item Регистрация шрифтов семейства Times New Roman (обычный, жирный, курсив, жирный курсив) для корректного отображения кириллицы. Библиотека \texttt{reportlab} не поддерживает кириллицу по умолчанию, поэтому требуется явная регистрация шрифтов с указанием пути к файлам \texttt{.ttf}.
	\item Настройка параметров страницы: формат A4, поля 15 мм со всех сторон.
	\item Формирование заголовочной части: название документа («Отчёт по проверке»), имя проверенного файла, дата и время проверки.
	\item Определение итогового статуса:
	\begin{itemize}
		\item если есть ошибки уровня \texttt{ERROR} — статус «документ содержит критические ошибки» (отображается красным цветом);
		\item если есть только предупреждения (\texttt{WARNING}) — «документ содержит предупреждения» (оранжевый);
		\item если ошибок нет — «документ соответствует требованиям» (зелёный).
	\end{itemize}
	\item Построение таблицы ошибок с колонками: номер ошибки, раздел проверки, тип (ошибка/предупреждение), местоположение в документе, описание проблемы.
	\item Стилизация таблицы: строки с ошибками выделяются красным фоном (\texttt{BackgroundColor = \#FEE2E2}), строки с предупреждениями — жёлтым (\texttt{BackgroundColor = \#FEF3C7}).
\end{enumerate}

Отчёт сохраняется в директорию \texttt{reports/} с именем, содержащим уникальный идентификатор, чтобы избежать конфликтов при одновременной работе нескольких пользователей.

\subsubsection{Генерация титульных листов}

Генерация титульных листов реализована через библиотеку \texttt{python-docx} без использования Microsoft Word. Это позволяет выполнять операцию быстро и без зависимости от установленного офисного пакета. В зависимости от типа документа используются разные наборы полей и шаблоны:

\begin{enumerate}
	\item Для методических указаний: название, дисциплина, направление подготовки, составитель, рецензенты (с возможностью добавления нескольких), УДК, описание.
	\item Для учебных пособий: автор, название, рецензенты, направления подготовки, значение А, ISBN, УДК, ББК, описание.
	\item Для монографии: авторы (с возможностью добавления нескольких), название, город, год, УДК, ББК, ISBN, описание.
\end{enumerate}

Алгоритм генерации единый для всех типов:

\begin{enumerate}
	\item Создаётся новый пустой документ DOCX.
	\item Устанавливаются поля страницы: верхнее 3 см, нижнее 2,25 см, левое 2,2 см, правое 2,3 см.
	\item Задаётся стиль Normal: шрифт Times New Roman, размер 16 pt.
	\item Последовательно добавляются абзацы с заданным выравниванием (центр, право, лево) и форматированием.
	\item Для каждого абзаца через вспомогательную функцию \texttt{add\_paragraph} устанавливаются параметры: текст, выравнивание, размер шрифта, жирное начертание, отступы до и после абзаца, межстрочный интервал.
	\item После заполнения первой страницы добавляется разрыв раздела (\texttt{doc.add\_section(WD\_SECTION\_START.NEW\_PAGE)}), и формируется вторая страница.
	\item Вторая страница содержит библиографическое описание, аннотацию и выходные данные — в зависимости от типа документа набор полей различается.
	\item Готовый документ сохраняется во временную директорию, после чего отправляется пользователю и удаляется с сервера.
\end{enumerate}

\subsubsection{Управление сессиями и обновление токенов}

Для поддержания авторизации пользователя без многократного ввода пароля используется двухуровневая система токенов: access-токен (срок действия 60 минут) и refresh-токен (срок действия 30 дней). Access-токен подписывается секретным ключом и содержит идентификатор пользователя и его email. Refresh-токен генерируется как случайная строка, его хеш сохраняется в базе данных вместе с метаинформацией (IP-адрес, User-Agent).

При получении запроса к защищённому ресурсу сервер проверяет access-токен: если он истёк, клиент получает ошибку \texttt{401 Unauthorized} и автоматически отправляет refresh-токен на эндпоинт \texttt{/auth/refresh}. Сервер проверяет, что refresh-токен не отозван (поле \texttt{revoked}), не истёк и принадлежит тому же пользователю, после чего выдаёт новый access-токен. Refresh-токен при этом остаётся без изменений, что позволяет обновлять access-токен многократно без повторного входа пользователя.

При выходе из системы соответствующий refresh-токен отмечается в базе данных как отозванный (\texttt{revoked = True}), что делает его непригодным для дальнейшего использования. Это позволяет администратору в случае подозрения на компрометацию учётной записи отозвать все токены пользователя одной командой.

\subsubsection{Взаимодействие с базой данных через ORM}

Для взаимодействия с PostgreSQL используется объектно-реляционное отображение SQLAlchemy. ORM позволяет работать с записями базы данных как с обычными объектами Python, скрывая детали формирования SQL-запросов. Каждая таблица описывается отдельным классом, наследующим от \texttt{Base}, с указанием типов полей, ограничений и связей.

Пример описания таблицы \texttt{manual}:

\begin{verbatim}
	class Manual(Base):
	__tablename__ = "manual"
	id_manual = Column(Integer, primary_key=True, index=True)
	manual_name = Column(String(255), nullable=False)
	fio_user = Column(String(255), nullable=False)
	faculty_code = Column(String(10), nullable=False)
	department_name = Column(String(255), nullable=False)
	file_hash = Column(String(64), unique=True, nullable=False)
	created_at = Column(DateTime, server_default=func.now())
\end{verbatim}

Сессии SQLAlchemy управляются через контекстный менеджер: при каждом запросе создаётся новая сессия, которая автоматически закрывается после окончания обработки. Это гарантирует, что соединения с базой данных не «утекают» и корректно возвращаются в пул.

Все запросы к базе данных используют параметризацию (через ORM-методы \texttt{filter} и \texttt{filter\_by}), что полностью исключает возможность SQL-инъекций. При сохранении документа с успешно пройденной проверкой хеш файла проверяется на уникальность; если запись с таким хешем уже существует, операция вставки не выполняется, и пользователь получает соответствующее уведомление.

\subsubsection{Предотвращение дублирования}

Для исключения повторной загрузки одинаковых документов используется вычисление хеш-суммы файла (алгоритм SHA-256). Хеш вычисляется блоками по 8192 байта, что позволяет обрабатывать файлы любого размера без полной загрузки в память.

При успешной проверке (документ не содержит ошибок) хеш сохраняется в базе данных вместе с метаинформацией: название документа, ФИО автора, код факультета, название кафедры, дата загрузки. При каждой новой загрузке система выполняет следующие действия:

\begin{enumerate}
	\item Вычисляется хеш загружаемого файла.
	\item Выполняется запрос к базе данных на поиск записи с таким же хешем.
	\item Если запись найдена, загрузка отклоняется, пользователю возвращается сообщение «Данный документ уже был загружен ранее».
	\item Если запись не найдена, документ сохраняется, и в базу добавляется новая запись.
\end{enumerate}

Это предотвращает накопление дубликатов и экономит дисковое пространство.

\subsubsection{Обработка ошибок и восстановление}

При работе с COM-интерфейсом Microsoft Word возможны различные сбои: документ может быть повреждён, Word может зависнуть или завершиться аварийно. В модуле реализована многоуровневая обработка исключений:

\begin{enumerate}
	\item На уровне COM-вызовов — каждый вызов метода Word обёрнут в блок \texttt{try-except}. При возникновении ошибки (например, попытка открыть повреждённый файл) выполняется корректное закрытие документа и завершение процесса Word через \texttt{self.word.Quit(False)}. Пользователь получает сообщение с указанием причины ошибки и предложением проверить файл.
	
	\item На уровне проверки документа — при сбое на любом этапе (фильтрация абзацев, извлечение параметров, группировка ошибок) проверка прерывается, пользователю возвращается код ошибки \texttt{500 Internal Server Error} с пояснением. Временные файлы удаляются в блоке \texttt{finally}, чтобы не засорять дисковое пространство.
	
	\item На уровне базы данных — при нарушении уникальности (попытка вставить дубликат) или других ограничениях целостности формируется понятное сообщение об ошибке без раскрытия деталей SQL-запроса.
\end{enumerate}

Дополнительно реализован механизм тайм-аутов: если проверка документа занимает более 600 секунд (10 минут), операция прерывается, и пользователь получает сообщение о необходимости уменьшить размер файла или разбить его на части.

\subsubsection{Хранение данных}

Для хранения метаданных используется таблица \texttt{manual} в PostgreSQL, содержащая следующие поля:

\begin{itemize}
	\item \texttt{id\_manual} — первичный ключ (автоинкремент);
	\item \texttt{manual\_name} — название документа;
	\item \texttt{fio\_user} — ФИО автора;
	\item \texttt{faculty\_code} — код факультета;
	\item \texttt{department\_name} — название кафедры;
	\item \texttt{file\_hash} — уникальный хеш (SHA-256);
	\item \texttt{created\_at} — дата и время загрузки.
\end{itemize}

Сами файлы хранятся в файловой системе в директориях \texttt{checked\_files/} (проверенные документы с выделенными ошибками) и \texttt{reports/} (PDF-отчёты). В базе хранятся только пути к этим файлам. Такой подход снижает нагрузку на СУБД и позволяет хранить файлы любого размера без ограничений.

\subsubsection{Логирование}

Все операции модуля логируются с использованием встроенного модуля \texttt{logging}. Настроены два уровня логирования:

\begin{enumerate}
	\item INFO — фиксируются штатные операции: загрузка файла, успешное завершение проверки, генерация титульного листа, отправка результата пользователю, авторизация и выход пользователя.
	
	\item ERROR — фиксируются сбои: невозможность открыть документ через COM-интерфейс, ошибка при сохранении файла, проблема с подключением к базе данных, тайм-аут проверки.
\end{enumerate}

Логи сохраняются в файл с ротацией по дате, что позволяет администратору системы анализировать причины отказов, отслеживать подозрительную активность и своевременно устранять неполадки.

\subsubsection{Заключение}

Разработанный модуль обработки учебно-методических документов реализует полный цикл работы с документами формата DOCX: приём файлов, автоматическую проверку оформления через COM-интерфейс Microsoft Word, группировку ошибок, формирование исправленной копии с визуальной разметкой, генерацию PDF-отчётов, создание титульных листов через \texttt{python-docx} (с различными шаблонами для методических указаний, учебных пособий и монографий), управление сессиями через двухуровневую систему токенов, предотвращение дублирования с помощью хеширования, надёжное хранение метаданных в PostgreSQL, а также многоуровневую обработку ошибок и логирование. Все компоненты модуля спроектированы с учётом требований к производительности, расширяемости и отказоустойчивости.

\subsection{База данных и её проектирование}

В качестве системы управления базами данных выбрана PostgreSQL. Эта СУБД обеспечивает высокую надёжность хранения информации, полную поддержку транзакций, соблюдение принципов ACID, а также располагает встроенными механизмами для контроля целостности данных на уровне ограничений и ссылочной целостности.

Подключение к базе данных реализовано через отдельный модуль database.py, который отвечает за конфигурацию параметров соединения, инициализацию подключений и их корректное завершение. В приложении задействован пул соединений — это позволяет переиспользовать уже установленные подключения, снижая накладные расходы на установку нового соединения для каждого запроса, и эффективно обслуживать большое количество одновременно работающих пользователей.

Для описания структур данных применяется объектно-реляционное отображение (ORM) SQLAlchemy. Каждая сущность предметной области представлена в коде отдельным классом, который описывает схему соответствующей таблицы: перечень полей, их типы, ограничения, а также связи с другими таблицами (один-ко-многим, многие-ко-многим). Такой подход упрощает взаимодействие с базой данных и снижает вероятность ошибок при формировании SQL-запросов.

\subsubsection{ER-модель данных}

На основе анализа предметной области были выделены основные сущности системы и связи между ними. ER-диаграмма отражает структуру данных и взаимосвязи между пользователями, кафедрами, факультетами и загруженными документами.

ER-диаграмма базы данных представлена на рисунке~\ref{fig:er_model}.

\begin{figure}[H]
	\centering
	\includegraphics[width=1\textwidth]{ERmodel.png}
	\caption{ER-диаграмма базы данных}
	\label{fig:er_model}
\end{figure}

В модели центральное место занимает сущность «manual», связанная с пользователями и организационной структурой (кафедры и факультеты). Также выделена сущность токенов обновления, обеспечивающая механизм аутентификации пользователей.

К числу основных сущностей базы данных относятся:

\begin{itemize}
	\item пользователи (users) — содержат учётные записи лиц, работающих с системой (авторы, преподаватели, сотрудники);
	\item факультеты (faculty) — подразделения университета, к которым привязаны пользователи и кафедры;
	\item кафедры (department) — структурные единицы внутри факультетов;
	\item документы (manual) — таблица, хранящая сведения о всех типах учебно-методических материалов (методические указания, учебные пособия, монографии), успешно прошедших проверку;
	\item токены обновления (refresh\_tokens) — используются для поддержания сеансов авторизации.
\end{itemize}

\subsubsection{Нормализация}

Структура базы данных спроектирована в соответствии с требованиями третьей нормальной формы.

Каждая таблица оснащена первичным ключом, все поля содержат неделимые значения, а взаимосвязи между сущностями организованы через механизм внешних ключей. Избыточность данных устранена путём декомпозиции информации на логически независимые таблицы (факультеты, кафедры, пользователи).

Помимо этого, в схеме предусмотрены дополнительные ограничения, гарантирующие целостность данных. К ним относятся уникальные индексы (например, для адреса электронной почты пользователя и хеш-суммы файла документа), а также проверочные условия (например, допустимые значения для роли пользователя).

\subsubsection{Реляционная модель данных}

На основе ER-модели построена реляционная модель базы данных, отражающая структуру таблиц и их взаимосвязи.

Схема реляционной модели представлена на рисунке~\ref{fig:relational_model}.

\begin{figure}[H]
	\centering
	\includegraphics[width=1\textwidth]{ReloModel.png}
	\caption{Реляционная модель базы данных}
	\label{fig:relational_model}
\end{figure}

Реляционная модель включает таблицы пользователей, факультетов, кафедр, документов (методических указаний, учебных пособий, монографий) и токенов обновления. Связи между таблицами реализованы через внешние ключи, что обеспечивает целостность данных и согласованность структуры.

Особое значение имеет таблица «manual», используемая для хранения информации о всех загружаемых документах и предотвращения их дублирования.

\subsubsection{Описание структуры базы данных}

Центральное место в базе данных занимает таблица «manual», предназначенная для хранения сведений о загружаемых документах независимо от их типа. Для каждой записи фиксируются следующие данные: название документа, фамилия, имя, отчество автора, код факультета, название кафедры, а также уникальная хеш-сумма файла. Хеш вычисляется на основе содержимого документа и служит для предотвращения дублирования — повторная загрузка одного и того же файла не приведёт к созданию лишней записи в базе.

Связь пользователей с организационной структурой обеспечивается через таблицы факультетов и кафедр. Каждый пользователь закреплён за конкретной кафедрой и факультетом, что позволяет отслеживать принадлежность авторов документов к соответствующим подразделениям и при необходимости формировать статистику по кафедрам или факультетам.

Файлы документов размещаются непосредственно в файловой системе сервера — это упрощает работу с большими объёмами данных и обеспечивает быстрый доступ к содержимому. В базе данных при этом сохраняются только метаданные (название, автор, дата загрузки) и уникальные идентификаторы, необходимые для поиска файла в файловом хранилище. Такое разделение позволяет снизить нагрузку на СУБД и эффективно масштабировать систему по мере увеличения количества загружаемых документов.

В таблицах~\ref{tab:faculty}--\ref{tab:tokens} приведено подробное описание структуры всех таблиц базы данных.

\begin{xltabular}{\textwidth}{|c|c|l|l|X|}
	\caption{Атрибуты таблицы «faculty» (Факультеты)\label{tab:faculty}}\\ \hline
	\centrow Key Type & \centrow Optionality & \centrow Column name & \centrow Data Type & \centrow Size \\ \hline
	\endfirsthead
	
	\continuecaption{Продолжение таблицы \ref{tab:faculty}}
	\centrow Key Type & \centrow Optionality & \centrow Column name & \centrow Data Type & \centrow Size \\ \hline
	\endhead
	
	PK & * & id\_faculty & SERIAL & \\ \hline
	& * & faculty\_name & VARCHAR & 255 \\ \hline
	UK & * & faculty\_code & INTEGER & \\ \hline
	& * & dean\_fio & VARCHAR & 255 \\ \hline
	&  & manual\_count & INTEGER & \\ \hline
\end{xltabular}

\begin{xltabular}{\textwidth}{|c|c|l|l|X|}
	\caption{Атрибуты таблицы «department» (Кафедры)\label{tab:department}}\\ \hline
	\centrow Key Type & \centrow Optionality & \centrow Column name & \centrow Data Type & \centrow Size \\ \hline
	\endfirsthead
	
	\continuecaption{Продолжение таблицы \ref{tab:department}}
	\centrow Key Type & \centrow Optionality & \centrow Column name & \centrow Data Type & \centrow Size \\ \hline
	\endhead
	
	PK & * & id\_department & SERIAL & \\ \hline
	& * & department\_name & VARCHAR & 255 \\ \hline
	UK & * & department\_code & INTEGER & \\ \hline
	&  & manual\_count & INTEGER & \\ \hline
	FK & * & id\_faculty & INTEGER & \\ \hline
\end{xltabular}

\begin{xltabular}{\textwidth}{|c|c|l|l|X|}
	\caption{Атрибуты таблицы «users» (Пользователи)\label{tab:users}}\\ \hline
	\centrow Key Type & \centrow Optionality & \centrow Column name & \centrow Data Type & \centrow Size \\ \hline
	\endfirsthead
	
	\continuecaption{Продолжение таблицы \ref{tab:users}}
	\centrow Key Type & \centrow Optionality & \centrow Column name & \centrow Data Type & \centrow Size \\ \hline
	\endhead
	
	PK & * & id\_user & SERIAL & \\ \hline
	& * & fio & VARCHAR & 255 \\ \hline
	UK & * & email & VARCHAR & 255 \\ \hline
	& * & password\_hash & VARCHAR & 255 \\ \hline
	& * & role & VARCHAR & 50 \\ \hline
	&  & created\_at & TIMESTAMP & \\ \hline
	FK & * & id\_faculty & INTEGER & \\ \hline
	FK & * & id\_department & INTEGER & \\ \hline
\end{xltabular}

\begin{xltabular}{\textwidth}{|c|c|l|l|X|}
	\caption{Атрибуты таблицы «manual» (Документы)\label{tab:manual}}\\ \hline
	\centrow Key Type & \centrow Optionality & \centrow Column name & \centrow Data Type & \centrow Size \\ \hline
	\endfirsthead
	
	\continuecaption{Продолжение таблицы \ref{tab:manual}}
	\centrow Key Type & \centrow Optionality & \centrow Column name & \centrow Data Type & \centrow Size \\ \hline
	\endhead
	
	PK & * & id\_manual & SERIAL & \\ \hline
	& * & manual\_name & VARCHAR & 255 \\ \hline
	& * & fio\_user & VARCHAR & 255 \\ \hline
	& * & department\_code & INTEGER & \\ \hline
	& * & faculty\_code & INTEGER & \\ \hline
	UK & * & file\_hash & VARCHAR & 64 \\ \hline
\end{xltabular}

\begin{xltabular}{\textwidth}{|c|c|l|l|X|}
	\caption{Атрибуты таблицы «refresh\_tokens» (Токены обновления)\label{tab:tokens}}\\ \hline
	\centrow Key Type & \centrow Optionality & \centrow Column name & \centrow Data Type & \centrow Size \\ \hline
	\endfirsthead
	
	\continuecaption{Продолжение таблицы \ref{tab:tokens}}
	\centrow Key Type & \centrow Optionality & \centrow Column name & \centrow Data Type & \centrow Size \\ \hline
	\endhead
	
	PK & * & id\_refresh\_token & SERIAL & \\ \hline
	UK & * & token\_hash & VARCHAR & 255 \\ \hline
	& * & expires\_at & TIMESTAMP & \\ \hline
	&  & created\_at & TIMESTAMP & \\ \hline
	& * & revoked & BOOLEAN & \\ \hline
	&  & user\_agent & VARCHAR & 500 \\ \hline
	&  & ip\_address & VARCHAR & 100 \\ \hline
	FK & * & id\_user & INTEGER & \\ \hline
\end{xltabular}